\section{Solving the interpolation equations}\label{sec:interpolation}

The equations \eqref{fourier:coefficient-equations} couple two coefficient
lists through the transformed part of the construction.  We solve them on
$\ell^1$, using the strict damping of the Fourier kernel to control every
row beyond a fixed finite set.

Let
\[
 X=\ell^1(\mathcal N;\RR^2),\qquad
 \norm{x}_1=\sum_{n\in\mathcal N}\bigl(\abs{c_n}+\abs{d_n}\bigr)
 \quad\text{for }x=((c_n,d_n))_{n\in\mathcal N}.
\]
When writing matrices, we order all $c$ coordinates first and all $d$
coordinates second, with increasing nodes within each group.  All operator
norms below are the induced $\ell^1$ norms, equivalently the maximum
absolute column sum for a finite matrix.  A subscript specifying a set of
nodes includes both coordinates at each node; a star denotes all positive
nodes.

For $x\in X$, let $p_x$ be the function \eqref{fourier:cardinal} with these
positive-node coefficients and with $C=c_0=d_0=0$.  Define the real linear
operator $S$ by
\begin{equation}\label{interpolation:S}
 (Sx)_{c,n}=-\frac{K_{p_x}(n)}{Q_n},\qquad
 (Sx)_{d,n}=\frac{D_nK_{p_x}(n)-K_{p_x}'(n)}{Q_n}.
\end{equation}
The bounds proved below show, in particular, that $S$ is a bounded operator
on $X$.

We use two successive decompositions of this operator. Set
\begin{equation}\label{interpolation:node-blocks}
 J=\{1,3,4\},\qquad E=\mathcal N\cap[7,100],\qquad
 I_f=J\cup E=\mathcal N\cap[1,100],\qquad D=S_{J,J}.
\end{equation}
A finite inverse handles the coordinates on $J$. A first Schur complement
then incorporates $E$, using separate bounds for its rows through $16$
and beyond $16$; a second incorporates the infinite tail beyond $100$.
The tail bounds apply to output rows uniformly over all input columns.
The first output rows require a separate, larger bound even for columns
indexed by large nodes.

\subsection{Uniform envelopes for the Fourier terms}

We use the following rational upper and lower bounds.  As throughout the
construction, every finite decimal denotes its exact rational value.
\begin{table}[htbp]
\centering
\begin{tabular}{c|rrr|rr}
$j$ & $b_j^*$ & $a_j^*$ & $g_j$ & $M_j$ & $\beta_j$\\ \hline
0 & 2.887 & 2.934 & 2.933 & --- & ---\\
1 & 2.665 & 2.713 & 2.440 & 56 & 0\\
2 & 2.218 & 2.269 & 1.567 & 48 & 4.8\\
3 & 1.804 & 1.859 & .900 & 35 & 3.17\\
4 & 1.486 & 1.548 & .482 & 25 & 1.946\\
5 & 1.250 & 1.320 & .224 & 19 & 1.217\\
6 & 1.073 & 1.151 & .0597 & 6.284 & .792
\end{tabular}
\caption{Envelopes for the folded Fourier integrals.}
\label{interpolation:envelopes}
\end{table}

At each endpoint $t_j=j/6$,
\[
 \abs{\lambda(t_j)}\le b_j^*,\qquad
 \abs{z(t_j)}\le a_j^*,\qquad
 \operatorname{Im}z(t_j)\ge g_j.
\]
On the interval $t_{j-1}\le t\le t_j$, the corresponding bounds are
\begin{equation}\label{interpolation:piece-envelopes}
 \abs{\lambda(t)}\le b_{j-1}^*,\quad
 \abs{z(t)}\le a_{j-1}^*,\quad
 \operatorname{Im}z(t)\ge g_j+\beta_j(t_j-t),\quad
 \abs{u_n(t)},\abs{v_n(t)}\le M_j.
\end{equation}
The density bounds hold uniformly for
$n\in\{0\}\cup\mathcal N$.

Here are direct ways to verify these estimates.  The identities
\[
 \abs{\lambda(t)}=\frac{1}{b\sqrt{t^2+h^2}},\qquad
 \abs{z(t)}^2=h^2+\frac{B^2-2Bh^2}{t^2+h^2},\qquad
 \operatorname{Im}z(t)=h\left(\frac{B}{t^2+h^2}-1\right)
\]
show that the first two quantities decrease on $[0,1]$.  The nonnegative
quantity $2hBt/(t^2+h^2)^2$, which is minus the derivative of
$\operatorname{Im}z$, has its minimum on each of the six intervals at an
endpoint.  Its only interior critical point on the positive axis is a
maximum.  These observations give the first three bounds in
\eqref{interpolation:piece-envelopes} by endpoint substitution.
For the densities, $e^{i\pi t_l n}$ depends only on $n$ modulo $12$,
whereas $\abs{e^{-i\pi tn}}=1$ on the real interval.  The norm of an
affine function of $t$ is convex.  Thus an upper bound for both folded
densities on the $j$th interval is
\begin{equation}\label{interpolation:density-check}
 2\max_{a\in A}\max\left\{
 \pi\abs{\sum_{l=j}^6P_le^{i\pi t_la}},\quad
 \pi^2\max_{u\in\{t_{j-1},t_j\}}
       \abs{\sum_{l=j}^6(t_l-u)P_le^{i\pi t_la}}
 \right\}.
\end{equation}
Substitution gives a value less than the displayed $M_j$ in every case.
For the constant input $C=1$, the absolute folded atom masses are
\[
 (m_0^*,\ldots,m_6^*)=(5,4,4,4,2,4,2).
\]

For $s\ge0$, set
\begin{equation}\label{interpolation:L}
 \begin{aligned}
 L_j(s)&=b_{j-1}^*e^{-\pi g_js}
       \min\left\{\frac16,\frac{1}{\pi\beta_js}\right\}
       &&(1\le j\le6),\\
 L_j^P(s)&=m_j^*b_j^*e^{-\pi g_js}
       &&(0\le j\le6).
 \end{aligned}
\end{equation}
When $\beta_js=0$, the minimum in the first line is understood to be
$1/6$.  Integrating the exponential envelope in
\eqref{interpolation:piece-envelopes} over a piece of length $1/6$ gives
its factor $L_j(s)$.  Differentiating the kernel contributes
$(i\pi z(t))^\ell$.  Consequently, for every integer $\ell\ge0$ and every
input $p$ of the form \eqref{fourier:cardinal},
\begin{equation}\label{interpolation:moment-bound}
 \abs{K_p^{(\ell)}(s)}
 \le w\sum_{j=1}^6 M_jL_j(s)(\pi a_{j-1}^*)^\ell
      +\abs C\sum_{j=0}^6L_j^P(s)(\pi a_j^*)^\ell,
 \qquad s\ge0,
\end{equation}
where $w=\sum_{n\in\{0\}\cup\mathcal N}(\abs{c_n}+\abs{d_n})$.
The argument applies first to finite lists and then to all such lists by
the finite-measure representation and absolute summability.

Both $Q_n$ and $D_n$ are $12$-periodic on the nodes.  Moreover,
\[
 L_j(n+12k)\le e^{-12k\pi g_j}L_j(n)\qquad(k\ge0),
\]
because the minimum factor in \eqref{interpolation:L} is nonincreasing.
For every real cutoff $x\ge0$, the value and derivative terms in
\eqref{interpolation:S}, together with
\eqref{interpolation:moment-bound}, therefore imply
\begin{equation}\label{interpolation:row-envelope}
 \norm{S_{>x,*}}
 \le
 \sum_{n\in\mathcal N\cap(x,x+12]}
 \sum_{j=1}^6
 \frac{1+\abs{D_n}+\pi a_{j-1}^*}{Q_n}\,
 \frac{M_jL_j(n)}{1-e^{-12\pi g_j}}.
\end{equation}
Here $>x$ means a strict inequality for the row node.  This is a bound
uniform over all input columns.  It applies as well to a single input coefficient
$c_0$ or $d_0$.

For completeness, this estimate also justifies the operator definition on
the whole space.  Apply the bound first to one positive-node basis column,
then sum with the absolute coefficients of a finitely supported list.
The case $x=0$ bounds the resulting output norm by a constant times the
input norm.  Finite lists are dense in $X$, so this map has a unique bounded
extension to $X$.  The basis measures converge in total variation by
Lemma~\ref{fourier:measure-lemma}; hence $K_{p_x}(n)$ and $K_{p_x}'(n)$
converge at each fixed node as the lists are truncated.  The extended map
therefore has exactly the coordinates in \eqref{interpolation:S}, not merely
the coordinates of an auxiliary matrix limit.

For $x=0,4,16,100$, respectively, the right side of
\eqref{interpolation:row-envelope} is less than
\[
 27.7,\qquad .272,\qquad .0145,\qquad 4.7\cdot10^{-10}.
\]
We shall use the slightly weaker bounds
\begin{equation}\label{interpolation:tails}
 \norm{S_{>0,*}}<28,\qquad
 \norm{S_{>4,*}}<.28,\qquad
 \norm{S_{>16,*}}<.015,\qquad
 \norm{S_{>100,*}}<5\cdot10^{-10}.
\end{equation}
For the input $C=1$, the analogue of
\eqref{interpolation:row-envelope} replaces $M_jL_j$ by $L_j^P$,
uses $a_j^*$ in place of $a_{j-1}^*$, and sums over $0\le j\le6$.
For rows greater than $100$, it gives the bound
\begin{equation}\label{interpolation:atom-tail}
 \sum_{n\in\mathcal N,\ n>100}
 \left(\frac{\abs{K_P(n)}}{Q_n}
 +\frac{\abs{D_nK_P(n)-K_P'(n)}}{Q_n}\right)
 <3.2\cdot10^{-8}.
\end{equation}
The scalar substitutions in these envelope estimates use only the explicit
constants and the node data.  Section~\ref{cert:scalar-bounds} records the
finite scalar bounds and the precise coverage of their supplied verifier.

\begin{remark}[Compactness of the interpolation operator]
\label{interpolation:compactness}
Let $P_M$ keep the output coordinates with node at most $M$.  Then $P_MS$
has finite rank.  The right side of \eqref{interpolation:row-envelope}
tends to zero as $x\to\infty$: there are only six nodes in any period,
the periodic $Q_n,D_n$ are uniformly bounded with $Q_n$ bounded away from
zero, and every $L_j(n)$ contains the factor $e^{-\pi g_jn}$ with $g_j>0$.
Thus $\norm{S-P_MS}\to0$.  The operator $S$ is a norm limit of finite-rank
operators and is compact on $X$.  This property describes the tail
structure; the explicit inverse below is obtained independently from the
quantitative block estimates.
\end{remark}

\subsection{The finite certificates}

There are $51$ nodes in $I_f$ and hence $102$ scalar coordinates.
The approximate lists $\widetilde x_1,\widetilde x_2$ are given by
Table~\ref{tab:coefficients}, with the indicated factor $10^{-10}$,
and are zero at positive nodes beyond $100$.  Append the fixed data
\eqref{fourier:fixed} to obtain $\widetilde p_i$, and write
$\widetilde K_i=K_{\widetilde p_i}$.

\begin{lemma}[Finite certificates]\label{interpolation:finite}
For each $\eta\in\{1,-1\}$, the matrix $I-\eta D$ is invertible, and
\begin{equation}\label{interpolation:finite-matrices}
 \norm{(I-\eta D)^{-1}}<32,\qquad
 \norm{(I-\eta D)^{-1}S_{J,E}}<3.7,\qquad
 \norm{S_{E\cap[7,16],J}}<.09.
\end{equation}
For $i=1,2$ and $n\in I_f$, the table satisfies
\begin{equation}\label{interpolation:finite-residuals}
 \begin{split}
 \abs{\widetilde c_{i,n}+\frac{\widetilde K_{3-i}(n)}{Q_n}}
 &<10^{-9},\\
 \abs{\widetilde d_{i,n}+\mathbf1_{\{i=1,\ n=1\}}
       +\frac{\widetilde K_{3-i}'(n)-D_n\widetilde K_{3-i}(n)}{Q_n}}
 &<10^{-9}.
 \end{split}
\end{equation}
For each $i$, the norm of the tabulated list, including node zero, is less
than $2.29$, and its norm on nodes greater than $40$ is less than $.00002$.
\end{lemma}

\begin{proof}
The primary interval calculation in Appendix~\ref{sec:certificates} proves
the exact bounds~\eqref{cert:matrix-table}, \eqref{cert:exterior-block}, and
\eqref{cert:residual-bounds}, using the quadrature assertion of
Lemma~\ref{cert:moment-error} for the analytic moment remainder.
The elementary argument in Section~\ref{cert:matrix-bounds} gives the
matrix conclusions, and~\eqref{cert:table-norms} gives the list norms.
\end{proof}

\subsection{Inversion on the whole coefficient space}

We first prove an inverse bound that keeps finite and tail residuals
separate.

\begin{lemma}\label{interpolation:inverse}
For each $\eta\in\{1,-1\}$, $I-\eta S$ has a bounded inverse on $X$.
If $r=r_f+r_t$, with $r_f$ supported on $I_f$ and $r_t$ on
$\mathcal N\cap(100,\infty)$, then
\begin{equation}\label{interpolation:inverse-bound}
 \norm{(I-\eta S)^{-1}r}_1
 \le91\norm{r_f}_1+2540\norm{r_t}_1.
\end{equation}
\end{lemma}

\begin{proof}
Fix $\eta$ and first restrict to the finite block $I_f$.  Put
\[
 V=(I-\eta D)^{-1},\qquad L=S_{E,J},\qquad U=S_{J,E}.
\]
The first block of tail estimates and the finite certificates give
\[
 \norm L
 \le\norm{S_{E\cap[7,16],J}}+\norm{S_{E\cap(16,100],J}}
 <.09+.015<.11.
\]
Eliminating the $J$ coordinates from
$(I-\eta S_{I_f,I_f})x_f=r_f$ leaves on $E$ the operator
\[
 I-\eta S_{E,E}-LVU.
\]
Its difference from the identity has norm less than
\[
 \norm{S_{E,E}}+\norm L\norm{VU}
 <.28+.11\cdot3.7=.687<1.
\]
It is therefore invertible by a Neumann series.  Solving this block and
substituting into the $J$ equations yields an inverse $V_f$ for the whole
finite compression, with
\begin{equation}\label{interpolation:finite-inverse-bound}
 \begin{split}
 \norm{V_f}
 &\le\max\left\{
 32+\frac{(1+3.7)\cdot.11\cdot32}{1-.687},\quad
 \frac{1+3.7}{1-.687}\right\}\\
 &<90.
 \end{split}
\end{equation}
Indeed, the two entries in the maximum bound the output norm arising
from a unit input on $J$ and on $E$, respectively.

Now write $X=X_f\oplus X_t$, where $X_t$ is the infinite tail space.
The full equation $(I-\eta S)x=r$ is equivalent, after elimination of
$x_f$, to
\begin{equation}\label{interpolation:tail-schur}
 \bigl(I-\eta S_{t,t}-S_{t,f}V_fS_{f,t}\bigr)x_t
 =r_t+\eta S_{t,f}V_fr_f.
\end{equation}
Put $\delta=5\cdot10^{-10}$.  By
\eqref{interpolation:tails},
\[
 \norm{S_{t,t}}<\delta,\qquad
 \norm{S_{t,f}}<\delta,\qquad
 \norm{S_{f,t}}<28.
\]
The operator subtracted from the identity in
\eqref{interpolation:tail-schur} thus has norm less than
\[
 \delta(1+90\cdot28)=2521\delta<1.
\]
Its Neumann series converges in the bounded operators on $X_t$.
Together with the finite inverse this proves invertibility on the full
space $X$, without any passage to a limit of finite-section solutions.
Quantitatively,
\[
 \norm{x_t}_1
 \le\frac{\norm{r_t}_1+90\delta\norm{r_f}_1}{1-2521\delta},
 \qquad
 \norm{x_f}_1\le90\norm{r_f}_1+2520\norm{x_t}_1.
\]
Adding these inequalities gives coefficients less than $91$ and $2540$
for $\norm{r_f}_1$ and $\norm{r_t}_1$, respectively, which proves
\eqref{interpolation:inverse-bound}.
\end{proof}

\begin{proposition}[Exact interpolation and coefficient enclosure]
\label{interpolation:solution}
There are unique real lists $x_1,x_2\in X$ satisfying
\eqref{fourier:coefficient-equations} with the fixed data
\eqref{fourier:fixed}.  For each $i=1,2$,
\begin{equation}\label{interpolation:coefficient-error}
 \sum_{n\in\mathcal N}
 \bigl(\abs{c_{i,n}-\widetilde c_{i,n}}
       +\abs{d_{i,n}-\widetilde d_{i,n}}\bigr)<.00003.
\end{equation}
The norm of each exact list, including its node-zero coefficients, is
less than $2.3$.  The corresponding functions satisfy all the
interpolation conditions \eqref{fourier:interpolation}.
\end{proposition}

The uniqueness here is within the specified coefficient space, with the
fixed node-zero data and constants. Triangular-shell values and first
derivatives do not determine an arbitrary radial Schwartz function and
its Fourier transform uniquely \cite[Theorem~1.7]{Sardari2021}.

\begin{proof}
Let $p_i^0$ denote \eqref{fourier:cardinal} with the fixed values
$(C_i,c_{i,0},d_{i,0})$ and all positive-node coefficients zero.  Let
$e\in X$ have its $d_1$ coordinate equal to one and all other coordinates
zero.  Define $g_i\in X$ by
\[
 (g_i)_{c,n}=-\frac{K_{p_{3-i}^0}(n)}{Q_n},\qquad
 (g_i)_{d,n}=
 \frac{D_nK_{p_{3-i}^0}(n)-K_{p_{3-i}^0}'(n)}{Q_n}
 -\mathbf1_{\{i=1\}}e_{d,n}.
\]
The estimates for the node-zero and constant columns show that these
are summable lists.  The equations to solve are precisely
\[
 x_1-Sx_2=g_1,\qquad x_2-Sx_1=g_2.
\]
Their sum is an equation for $I-S$, and their difference is an equation
for $I+S$.  Lemma~\ref{interpolation:inverse} gives unique real solutions
and hence unique $x_1,x_2$.

It remains to locate these solutions relative to the table.  Set
\[
 \rho_i=\widetilde x_i-S\widetilde x_{3-i}-g_i.
\]
There are $102$ scalar coordinates in $I_f$, so
\eqref{interpolation:finite-residuals} gives
\[
 \norm{(\rho_i)_f}_1<1.1\cdot10^{-7}=:R_f.
\]
On the tail, both the tabulated positive-node list and the vector $e$
vanish.  The density-column estimate, now including node zero, and
\eqref{interpolation:atom-tail} imply
\[
 \norm{(\rho_i)_t}_1
 <2.29\cdot5\cdot10^{-10}+.017\cdot3.2\cdot10^{-8}=:R_t.
\]
Here $2.29$ bounds the whole tabulated coefficient list, and $.017$
bounds both fixed constants in absolute value.

The errors $x_i-\widetilde x_i$ are recovered by taking half the sum
and half the difference of the solutions with residuals
$-(\rho_1+\rho_2)$ and $-(\rho_1-\rho_2)$.  Applying
\eqref{interpolation:inverse-bound} to both gives, for each list,
\[
 \norm{x_i-\widetilde x_i}_1
 \le2(91R_f+2540R_t)
 =.00002860012<.00003.
\]
This proves \eqref{interpolation:coefficient-error}.  Since the fixed
node-zero coefficients agree, the total exact coefficient norm is
less than $2.29+.00003<2.3$.  Finally,
\eqref{fourier:coefficient-equations} is equivalent to
\eqref{fourier:interpolation} by Lemma~\ref{fourier:equivalence}.
\end{proof}
